\chapter{Anger Management}

\section*{Definition and Overview}
Anger is a normal human emotion that serves as a signal that something is wrong or that a boundary has been crossed. Anger becomes a problem when it is frequent, intense, disproportionate to the situation, or expressed through aggression or violence, harming relationships, work, and physical and mental health. Anger management refers to the skills, techniques, and structured programmes that help people recognise early warning signs, understand their triggers, and respond in calmer and more constructive ways.

Problem anger is a treatable condition rather than a fixed personality trait. With education, psychological therapy, and lifestyle change, most people can learn to experience anger without letting it control their behaviour. Anger management is relevant across the lifespan, from children and adolescents with frequent temper outbursts to adults and older people whose irritability is straining their relationships or health.

\section*{Epidemiology}
Anger that is frequent, intense, or poorly controlled affects a substantial minority of the population. Community surveys suggest that around one in five people report difficulty controlling their temper, and that a smaller proportion exhibit recurrent physical or verbal aggression. Men are more likely to display physical aggression, whereas women report comparable levels of internal anger experience and verbal outbursts.

Problem anger is more common in younger adults, in people under chronic stress, and in those with depression, anxiety, post-traumatic stress disorder (PTSD), attention deficit hyperactivity disorder (ADHD), or alcohol and drug problems. People who grew up in households where aggressive behaviour was modelled, and those who have experienced trauma or maltreatment, are at higher risk. Anger problems frequently coexist with other mental health difficulties, and they are over-represented in people seeking help for relationship conflict and in forensic populations.

\section*{Aetiology and Pathophysiology}
Anger arises from a complex interaction of biological, psychological, and social factors:

\begin{itemize}
    \item \textbf{Biological:} the sympathetic nervous system and stress response release adrenaline, cortisol, and other hormones, producing a state of physiological arousal that readies the body for fight or flight
    \item \textbf{Genetics and temperament:} some individuals are constitutionally quicker to experience strong anger, and there is evidence of heritable components to trait anger
    \item \textbf{Learned behaviour:} anger is often modelled on the responses of parents, peers, and cultural surroundings during development
    \item \textbf{Psychological factors:} rigid expectations, a sense of being disrespected, low frustration tolerance, and distorted thinking (such as personalising events or catastrophising)
    \item \textbf{Stress and fatigue:} poor sleep, work overload, chronic pain, and life pressures lower the threshold at which anger is triggered
    \item \textbf{Substance use:} alcohol and stimulant drugs reduce inhibition and increase the risk of aggressive outbursts
    \item \textbf{Mental health conditions:} depression, anxiety, PTSD, and ADHD frequently present with irritability and anger as prominent symptoms
\end{itemize}

The neural circuits involved include the amygdala, which detects threats and generates emotional responses, and the prefrontal cortex, which normally exerts inhibitory control over these responses. When prefrontal regulation is weakened by stress, fatigue, intoxication, or underlying mental illness, angry reactions are more easily triggered and harder to contain.

\section*{Clinical Features}
Anger manifests across several domains:

\begin{itemize}
    \item \textbf{Physical signs:} racing heart, rapid breathing, clenched jaw or fists, muscle tension, sweating, and a hot or flushed feeling
    \item \textbf{Emotional signs:} irritability, feeling provoked or resentful, a sense of rising frustration, and feeling unable to stop the emotion
    \item \textbf{Behavioural signs:} shouting, arguing, slamming doors, throwing objects, and physical aggression
    \item \textbf{Verbal and psychological aggression:} insults, threats, intimidation, and attempts to control or frighten others
    \item \textbf{Withdrawal:} silent resentment, sulking, and the cold shoulder as a way of punishing others
    \item \textbf{Consequences:} relationship conflict, workplace problems, driving aggression, and feelings of guilt or shame after outbursts
\end{itemize}

The pattern is typically described as a build-up phase of rising tension, an explosion or expression phase, and a post-episode phase of regret, exhaustion, or residual resentment.

\section*{Differential Diagnosis}
Problem anger is often a symptom rather than a diagnosis in itself, so the following should be considered:

\begin{itemize}
    \item \textbf{Intermittent explosive disorder:} recurrent, impulsive aggressive outbursts that are grossly out of proportion to the trigger
    \item \textbf{Depression, anxiety, or PTSD:} conditions in which irritability and anger are common features
    \item \textbf{ADHD:} poor impulse control and emotional dysregulation
    \item \textbf{Alcohol and substance use disorders:} intoxication, withdrawal, and dependence all lower the threshold for aggression
    \item \textbf{Personality factors:} traits of borderline, antisocial, or narcissistic personality that impair emotional regulation
    \item \textbf{Medical causes:} chronic pain, sleep deprivation, thyroid dysfunction, and cognitive decline or dementia in older people
    \item \textbf{Normal grief and adjustment:} situational anger in response to loss or major life change
\end{itemize}

A thorough assessment distinguishes primary anger problems from anger that is driven by an underlying treatable condition.

\section*{When to Seek Medical Care}
People should seek help from a general practitioner or psychologist if anger is frequent, intense, or out of proportion to situations; if it is damaging relationships, work, or health; if they feel unable to control their temper; or if family, friends, or employers have raised concerns. Children and teenagers with persistent, escalating anger or aggression should be assessed early, as should older people with a new change in irritability.

\textbf{Call 000 (or local emergency services) for:}
\begin{itemize}
    \item Thoughts of harming yourself or someone else
    \item A violent situation that is escalating or out of control
    \item Feeling that you may act on an urge to hurt another person
    \item Someone who is at immediate risk because of your anger or theirs
\end{itemize}

\section*{Diagnosis and Investigations}
Assessment of anger problems is based mainly on clinical interview and structured measures:

\begin{itemize}
    \item \textbf{Clinical interview:} detailed history of the anger -- its frequency, intensity, triggers, typical pattern of escalation, and consequences for relationships, work, and health
    \item \textbf{Mental health assessment:} screening for depression, anxiety, trauma, ADHD, and substance use, since these frequently drive or worsen anger
    \item \textbf{Standardised questionnaires:} validated anger scales (for example the State--Trait Anger Expression Inventory) to quantify the problem and monitor change
    \item \textbf{Risk assessment:} direct enquiry about any history of violence, self-harm, threats, or harm to others, including domestic and family violence
    \item \textbf{Physical examination and investigations:} targeted examination and blood tests if the history suggests a medical cause such as thyroid disease or chronic pain
    \item \textbf{Collateral information:} where appropriate, information from a partner or family member to understand the impact of the anger
\end{itemize}

\section*{Management}
Effective management combines structured psychological therapy with lifestyle and, where needed, medical treatment:

\begin{itemize}
    \item \textbf{Cognitive behaviour therapy (CBT):} the best-evidenced approach; it teaches people to identify automatic thoughts, challenge unhelpful beliefs, and use calming self-talk
    \item \textbf{Structured anger management programmes:} group or individual courses covering relaxation, problem-solving, communication, and relapse prevention
    \item \textbf{Relaxation training:} slow breathing, progressive muscle relaxation, and mindfulness to dampen physiological arousal
    \item \textbf{Communication skills:} assertiveness training, expressing feelings calmly with "I" statements, and active listening instead of blaming
    \item \textbf{Time-out strategies:} deliberately stepping away from a heated situation until arousal subsides
    \item \textbf{Treating underlying conditions:} psychological therapy or medication for depression, anxiety, trauma, or ADHD, and reduction of alcohol and drug use
    \item \textbf{Lifestyle measures:} regular exercise, adequate sleep, and stress management to raise the threshold for anger
    \item \textbf{Supporting services:} for people at risk of violence, referral to specialist domestic violence or mental health services is essential
\end{itemize}

\section*{Complications}
\begin{itemize}
    \item Damage to intimate and family relationships, including children who witness aggression
    \item Workplace problems, disciplinary action, job loss, and financial stress
    \item Physical and verbal aggression, with risk of injury to others or to oneself
    \item Legal consequences, including assault charges, police involvement, and protective orders
    \item Domestic and family violence, with serious physical and psychological harm to victims
    \item Long-term health effects of chronic stress, including hypertension and cardiovascular disease
    \item Depression, anxiety, substance misuse, and social isolation
\end{itemize}

\section*{Prognosis and Outcomes}
Problem anger is treatable, and most people who complete structured anger management programmes report fewer, less intense episodes, better control, and improved relationships. CBT-based programmes produce meaningful improvement in roughly two-thirds of participants, and gains are often maintained at follow-up. Outcomes are better when the person recognises the problem early, remains engaged with treatment, and simultaneously addresses alcohol use, stress, sleep, and any underlying mood or trauma-related difficulties. Without help, problematic anger frequently persists and can escalate over time, with worsening relationship and health consequences.

\section*{Prevention and Self-Care}
\begin{itemize}
    \item Learn to recognise your early warning signs of rising anger and act before escalation begins
    \item Practise slow, deep breathing whenever you notice tension building
    \item Use time-out: leave the situation, allow your body to calm down, and return to discuss the issue later
    \item Identify your common triggers and plan in advance how to handle them
    \item Protect your sleep, exercise regularly, and limit alcohol and other drugs
    \item Practise expressing feelings calmly and assertively rather than aggressively
    \item Challenge unhelpful thoughts such as assuming others intend to provoke you
    \item Seek help early from a GP, psychologist, or structured anger management course
    \item If you sense violence may occur, leave the situation and get support immediately
\end{itemize}

\section*{Sources and Further Reading}
The following sources provide reliable further information on anger and anger management:

\begin{itemize}
    \item NHS. \textit{Overview of anger problems, triggers, and self-help and treatment options.} https://www.nhs.uk/mental-health/conditions/anger/
    \item American Psychological Association. \textit{Resources on controlling and managing anger.} https://www.apa.org/topics/anger
    \item Mayo Clinic. \textit{Information on anger management techniques and treatment.} https://www.mayoclinic.org/healthy-lifestyle/adult-health/in-depth/anger-management/art-20045434
    \item Mind (UK). \textit{Understanding and managing anger, including self-help approaches.} https://www.mind.org.uk/information-support/types-of-mental-health-problems/anger/
    \item NICE. \textit{Guidance relevant to antisocial behaviour and anger in young people.} https://www.nice.org.uk/guidance/ng58
\end{itemize}
